\documentclass[11pt]{article}

% --- packages ---------------------------------------------------------------
\usepackage[utf8]{inputenc}
\usepackage[T1]{fontenc}
\usepackage[margin=1in]{geometry}
\usepackage{microtype}
\usepackage{amsmath}
\usepackage{booktabs}
\usepackage{graphicx}
\usepackage{xcolor}
\usepackage{listings}
\usepackage[numbers,sort&compress]{natbib}
\usepackage[colorlinks=true,allcolors=blue]{hyperref}

% --- listings (JSON / ascii diagrams) ---------------------------------------
\definecolor{codebg}{gray}{0.96}
\lstset{
  basicstyle=\ttfamily\footnotesize,
  backgroundcolor=\color{codebg},
  frame=single,
  framesep=4pt,
  rulecolor=\color{gray!40},
  breaklines=true,
  columns=fullflexible,
  keepspaces=true,
  showstringspaces=false,
  aboveskip=1em,
  belowskip=1em,
}

\newcommand{\code}[1]{\texttt{\small #1}}

% --- title ------------------------------------------------------------------
\title{\textbf{DSLs as Harnesses:}\\[2pt]
\large Declarative Applications-as-Data as a Reliability Substrate for LLM Agents}

\author{
  Satoshi Nakajima\\
  \small The Singularity Society\\
  \small \texttt{satoshi.nakajima@gmail.com}
}

\date{\today}

\begin{document}
\maketitle

\begin{abstract}
A recurring lesson of applied agent engineering is that the \emph{harness}---the
designed environment an agent operates in---governs reliability more than the
underlying model. We argue that a \emph{domain-specific language} (DSL) is one of
the most effective harnesses available for autonomous, long-horizon agents,
precisely because it gives up the unbounded expressiveness of general-purpose
code in exchange for a small, legible, checkable surface. We make this concrete
by describing \textbf{Collections}, a declarative \emph{applications-as-data}
system implemented in MulmoClaude, an open-source LLM-agent application. In
Collections, a single JSON schema defines an application's data model,
relationships, computed values, recurrence, user interface, external-data
acquisition, and workflows; record files are the database; and the language model
is the runtime that operates inside the schema rather than writing imperative
code. We show how four well-known harness properties---schema-as-system-of-record,
a validator as a free feedback loop, a grammar as a forcing function, and a unit
of the language as a unit of context---fall out of this single design decision,
and how a deliberate \emph{reliability dial} divides what the host guarantees
deterministically (schema validity, derived-value computation, recurrence math,
sandboxed rendering, scheduled ingestion) from what the model supplies (judgment,
language, domain modeling). We report a longitudinal single-user deployment in
which 26 applications, 6 bespoke sandboxed UI views, and 3 external feeds accreted
over roughly two months without any host code being written per application. We
discuss the design's principal hazard---the escape hatch from the declarative
surface to arbitrary code---and its limitations, including non-replayable
execution and the single-user, single-system scope of our evidence.
\end{abstract}

\noindent\textbf{Keywords:} LLM agents; agent harness; domain-specific languages;
end-user software engineering; low-code; applications as data; declarative systems.

% ===========================================================================
\section{Introduction}
\label{sec:intro}

The prevailing lesson of applied AI engineering in 2024--2026 is that the
\emph{harness}---the designed environment an agent operates inside: the tools it
can call, the format of the information it receives, the feedback that catches its
mistakes, and the scaffolding that lets it hand work to its future self---matters
more than the model alone~\citep{yang2024sweagent,yao2023react}. The same model,
given a better harness, produces dramatically better work. For a language model,
the interface is not a convenience layer wrapped around the intelligence; it is
the cognitive architecture.

This paper argues a narrower and, we believe, more actionable point: \textbf{a
domain-specific language (DSL) is one of the most powerful harnesses one can give
an agent, precisely because it gives up power.} A general-purpose programming
language is Turing-complete; that is its glory and, for an autonomous agent, its
hazard. There are infinitely many ways to express any behavior, no two codebases
agree on conventions, and nothing in the language tells the agent whether a given
program is \emph{correct for this domain}. Freedom under uncertainty is where
agents thrash: they explore broadly, drift toward locally plausible patterns, and
accumulate subtle errors that surface much later. A DSL inverts every one of those
properties---it expresses a bounded set of things, in a bounded number of ways,
and typically ships with a parser and validator that can decide well-formedness
before anything executes.

We ground the argument in a working system. \textbf{Collections} is a feature of
MulmoClaude,\footnote{\url{https://github.com/receptron/mulmoclaude}} an
open-source agent application built on the Claude Agent SDK whose
founding philosophy is \emph{the workspace is the database; files are the source
of truth; the model is the intelligent interface.} A collection is defined by a
single JSON \emph{schema} that simultaneously specifies a data model, inter-record
relationships, spreadsheet-like derived computations, completion and notification
rules, host-driven recurrence, a rendered user interface, optional bespoke HTML
views, and optional scheduled ingestion of external data. The records are plain
JSON files on disk. The host platform contains \emph{zero} domain knowledge; it
understands only generic concepts (fields, relationships, derived values, views,
actions, ingestion). Everything domain-specific lives in the schema, and the
language model operates the application by reading and writing records under the
schema's constraints.

The significance for reliability is that a self-built, model-operated application
need not be ``flaky chat with extra steps.'' The DSL is a contract that pins down
what must be deterministic while the model supplies flexibility and judgment, and
the two live on opposite sides of a deliberate line we call the \emph{reliability
dial} (\S\ref{sec:dial}). A schema \emph{validator} makes the boundary
self-enforcing: a structurally incoherent application is rejected at load, with an
actionable message, before it can run.

\paragraph{Contributions.}
\begin{enumerate}
  \item We articulate \textbf{DSL-as-harness}: a mapping from a single design
  decision (constrain the agent's authoring surface to a validated DSL) to four
  established harness properties, and an argument for why the case strengthens,
  rather than weakens, as models improve (\S\ref{sec:harness}).
  \item We describe \textbf{Collections}, a concrete declarative
  applications-as-data system, including its schema DSL, relationships, derived
  formulas, host-driven recurrence, sandboxed custom views, and declarative feed
  ingestion (\S\ref{sec:collections}).
  \item We formalize the \textbf{host/model responsibility split} and the
  \emph{reliability dial}, and explain the validator and escape-hatch mechanisms
  that keep the surface both safe and expressive (\S\ref{sec:dial},
  \S\ref{sec:escape}).
  \item We report a \textbf{longitudinal single-user deployment}: 26 applications,
  6 bespoke sandboxed views, and 3 external feeds accreted over $\sim$2 months
  with no per-application host code (\S\ref{sec:casestudy}).
  \item We give an \textbf{honest account of the limitations}: the escape-hatch
  hazard, non-replayable execution, and the scope of our evidence
  (\S\ref{sec:limitations}).
\end{enumerate}

% ===========================================================================
\section{Background and Related Work}
\label{sec:related}

\paragraph{Agent harnesses and agent--computer interfaces.}
\citet{yang2024sweagent} showed that a purpose-built agent--computer interface
(ACI)---structured tools, capped and guided search, and a linter that rejects
syntactically broken edits at the moment of introduction---improves software-engineering
agent performance substantially with a fixed model. \citet{yao2023react}
interleave reasoning and acting; \citet{schick2023toolformer} teach models to call
external tools. These works establish the empirical primacy of the surrounding
environment. Our contribution is to identify the \emph{DSL} as a particularly
sharp instance of such an environment and to instantiate it as an
applications-as-data substrate.

\paragraph{Self-improving and memory-augmented agents.}
Voyager~\citep{wang2023voyager} accretes a reusable skill library, and Generative
Agents~\citep{park2023generative} maintain a memory stream they reflect over. Both
let an agent extend itself over time. Collections differ in substrate: the agent's
accreted capabilities are \emph{declarative DSL artifacts on disk}, not learned
weights or imperative code, which makes them inspectable, diffable, and validated.
A complementary line of work makes the event log the primary
artifact---\citet{ynakajima2026log} prioritizes an append-only log to obtain
deterministic replay and cheap forking. We deliberately forgo byte-for-byte replay
(the model step is nondeterministic) in exchange for an agent whose configuration
is plain, readable, validated files; \S\ref{sec:limitations} returns to this
trade.

\paragraph{Domain-specific languages.}
The software-engineering literature has long studied when and how a deliberately
limited language pays off~\citep{mernik2005dsl,fowler2010dsl}: a DSL trades general
expressiveness for analyzability, a smaller error surface, and tooling (parsers,
validators) that come ``for free.'' We apply this classical argument to a new
consumer---an autonomous language-model agent---and observe that the properties
DSLs were always praised for map directly onto modern harness-engineering
guidance.

\paragraph{End-user software engineering and low-code.}
End-user programming has sought for decades to let non-programmers build software,
from the spreadsheet onward~\citep{nardi1993small,ko2011enduser}; commercial
low-code/no-code platforms (Airtable, Notion, Retool, PowerApps) continue this
lineage. The persistent constraint in that tradition is that the user must work
inside a constrained authoring surface \emph{and} perform the assembly. Collections
keeps the constrained, validated surface but moves authorship to the agent: the
user expresses intent in natural language and the agent emits the DSL. We position
this as a shift from \emph{no-code} (a human composes blocks) to
\emph{intent-driven} (an agent composes the DSL program from intent).

% ===========================================================================
\section{A DSL Is a Harness}
\label{sec:harness}

When you hand an agent a DSL you are not merely giving it a tool; you are reshaping
the space it must search. That reshaping makes four classic harness patterns fall
out of a single design decision.

\paragraph{(1) The schema is the system of record.}
Harness engineering insists that anything the agent cannot read in context
effectively does not exist; intent must live in a machine-readable artifact. A DSL
satisfies this by construction: its schema \emph{is} the specification. When the
agent writes a DSL document, the user's intent is captured as a structured,
inspectable object rather than dissolved into the incidental details of imperative
code. It can be diffed, versioned, validated, and reasoned about without running
it.

\paragraph{(2) The validator is a free feedback loop.}
The highest-leverage component in the SWE-agent study was a linter that rejected
broken edits at the moment of introduction, before the error could propagate
through later steps~\citep{yang2024sweagent}. Every DSL ships with this loop:
the instant the agent emits a malformed structure, a parser or schema checker
rejects it with a localized, actionable message. The feedback loop came bundled
with the language.

\paragraph{(3) The grammar is a forcing function.}
Because the vocabulary is finite, the room to hallucinate is structurally smaller.
The agent cannot invent an option absent from the grammar; it can only compose the
primitives the language provides. Constraint, here, is a feature.

\paragraph{(4) The unit of the language is the unit of context.}
A well-designed DSL has a natural decomposition---a field, a record, a rule---which
becomes the agent's editing unit. It can work on one piece without holding the
whole document in mind, and the language's structure, not an ad-hoc heuristic,
defines the seams (progressive disclosure).

\paragraph{Why this matters more as models improve.}
A tempting intuition holds that DSLs are a crutch for weak models. We argue the
opposite. The case for a constrained, validated surface is strongest exactly where
we want to delegate \emph{autonomously and over long horizons}: where no human
reviews each step, where output must be verifiable, and where execution must be
deterministic and reversible. A stronger model inside a DSL harness is not wasted
capability; it is capability aimed at the part of the problem that genuinely
requires judgment---what to model, what to express---while the language guarantees
the rest.

% ===========================================================================
\section{The Collections System}
\label{sec:collections}

Collections applies the DSL-as-harness idea to software itself. A collection is:
\begin{center}
\code{schema.json} \;+\; \code{records/*.json} \;+\; model \;=\; application.
\end{center}
There is no database server, no migration framework, and no application-specific
backend. The host understands only generic capabilities; everything
domain-specific lives in the collection.

\subsection{The schema DSL}
A schema declares typed \emph{fields}. The current implementation supports 17
field types, including \code{string}, \code{text}, \code{number}, \code{date},
\code{datetime}, \code{boolean}, \code{markdown}, \code{enum}, \code{money},
\code{image}, \code{file}, a nested \code{table} of sub-records, two relationship
types (\code{ref}, \code{embed}), a computed \code{derived} field, and a
\code{toggle} that projects an enum. Relationships are first-class: a \code{ref}
stores a foreign key and the host renders pickers, links, and navigation
generically; an \code{embed} pulls a fixed record from another collection for
display.

\subsection{Computation without code}
A \code{derived} field behaves like a spreadsheet formula, but formulas may follow
references into other collections:
\begin{lstlisting}
"total": {
  "type": "derived", "display": "money", "currencyField": "currency",
  "formula": "sum(lineItems[].quantity * lineItems[].rate) * (1 + taxRate)"
}
\end{lstlisting}
The host evaluates derived fields, and crucially returns them \emph{to the model}
as well as to the screen, so the agent reasons over the same numbers the user
sees. A portfolio can revalue itself when a quote changes elsewhere, with no
synchronization code.

\subsection{Completion, notification, and host-driven recurrence}
A schema can mark which field and values count as ``done''
(\code{completionField} / \code{completionDoneValues}), gate notifications on a
date (\code{triggerField} / \code{triggerLeadDays}), and---via a \code{spawn}
block---provision the next instance of a recurring obligation on a civil-date
cadence when the current one closes:
\begin{lstlisting}
"triggerField": "dueOn", "triggerLeadDays": 3,
"spawn": {
  "every": { "unit": "month", "interval": 1, "dayOfMonth": 10 },
  "carry": ["payee", "amount"], "set": { "status": "pending" }
}
\end{lstlisting}
The author declares \emph{what recurs and when to nudge}; the host's reconciler
owns the deterministic half---raising and clearing reminders and advancing each
cycle with create-if-absent (idempotent) semantics. No timer, cron expression, or
notification code is written.

\subsection{Custom views: bespoke UI as data}
Host-rendered field types give every collection a list, a detail form, and---when
declared---calendar and Kanban layouts for free. When field types cannot present
the data well (a portfolio as an allocation chart, a set of places as a map, a
word list as flashcards), a collection may carry a \emph{custom view}: an HTML
file the agent authors and the host serves.
\begin{lstlisting}
"views": [
  { "id": "allocation", "label": "Allocation",
    "file": "views/allocation.html", "capabilities": ["read", "write"] }
]
\end{lstlisting}
A view is an ordinary web page (it may load a charting library from a curated CDN
allow-list) but it is rendered inside a sandboxed \code{iframe} under a strict
Content-Security-Policy: an opaque origin, \code{default-src 'none'}, and network
access only back to the host. It receives its records through a capability-scoped,
time-limited token rather than ambient access; the declared \code{capabilities}
(\code{read}, \code{write}) bound what that token permits. A custom view is thus
not a hole in the harness but the same bargain as the rest of the system: the
model supplies the bespoke surface, the host supplies the deterministic, sandboxed
boundary, and the presentation layer too becomes inspectable, versioned data.

\subsection{Feeds: acquisition as data}
Some applications track an external source---a podcast's episodes, a newspaper's
headlines, a JSON API. Such a collection declares an \code{ingest} block and
becomes a \emph{feed}:
\begin{lstlisting}
"ingest": {
  "kind": "rss", "url": "https://lexfridman.com/feed/podcast/",
  "schedule": "daily", "idFrom": "guid",
  "map": { "headline": "title", "url": "link", "published": "pubDate" },
  "maxItems": 100
}
\end{lstlisting}
\code{kind} selects a retriever (\code{rss}, \code{atom}, \code{http-json});
\code{map} projects source fields onto collection fields; \code{schedule}
(\code{hourly}/\code{daily}/\code{weekly}/\code{on-demand}) sets cadence;
\code{maxItems} caps retention. The host owns fetching, parsing, field mapping,
de-duplication by id, pruning, and cursor/failure tracking. Because a feed is
still a collection, the user can layer their own fields onto the ingested
ones---a rating, a note, a resume position---turning a passive mirror into a
read--write application over an external source.

% ===========================================================================
\section{The Reliability Dial}
\label{sec:dial}

The reason a self-built, model-operated application is reliable rather than flaky
is a single principle: \textbf{the DSL is a contract that pins down what must be
deterministic, while the model supplies flexibility, judgment, and language.}
Table~\ref{tab:split} summarizes the split.

\begin{table}[t]
\centering
\small
\begin{tabular}{@{}ll@{}}
\toprule
\textbf{Host guarantees (deterministic)} & \textbf{Model owns (judgment)} \\
\midrule
Schema validity (load-time)        & What to model \\
Path safety, record storage        & What a record means \\
Derived-value computation          & How a workflow should proceed \\
Recurrence / civil-date math       & When to hand off to a human \\
Sandboxed view rendering (CSP)     & How to present data (authoring the view) \\
Scheduled ingestion + de-dup       & Which external source, mapped how \\
\bottomrule
\end{tabular}
\caption{The host/model responsibility split. Everything on the left runs
identically every time; everything on the right is where the model's flexibility
is wanted.}
\label{tab:split}
\end{table}

A schema validator makes the boundary self-enforcing. A structurally incoherent
application is rejected at load with an actionable reason---examples include a
money field with neither a literal currency nor a per-record currency field, a
\code{ref} pointing at a non-existent collection, a \code{triggerField} that does
not name a real date field, and a \code{spawn} whose successor would be born
already matching its own predicate (an unbounded respawn).

The craft is \emph{where the dial sits}. Push everything into the DSL and you
rebuild the brittle ``can't-express-this'' cage that made no-code platforms
frustrating; push everything to the model and you lose the durable contract that
made it an application at all. The guiding rule is to \emph{extend the declarative
layer only when it outperforms the agent}.

% ===========================================================================
\section{The Escape Hatch}
\label{sec:escape}

The same rigidity that makes a DSL a good harness is its failure mode: when the
language cannot express what the user needs, an agent will either give up or
contort the DSL into something it was never meant to hold. Every well-designed DSL
answers this with a deliberate exit to a more expressive layer.

Collections provides two. First, a per-record or collection-level \code{action}
hands off to a seeded chat---a full-tool agent---exactly when a task needs human
judgment the form cannot capture (``generate the invoice PDF,'' ``what was the
amount, and was it actually paid?''). This is business logic as prose, behind a
declarative door. Second, the \emph{custom view} (\S\ref{sec:collections}) is the
escape for presentation: when host-rendered field types cannot express a layout,
the schema drops to arbitrary HTML in a sandbox rather than distorting the
field-type vocabulary. The art is not to maximize constraint but to choose
\emph{where} the constraint binds and to provide a clean, legible path out where it
does not. A harness with no escape hatch eventually forces the agent to fail or to
lie, and both are worse than a slightly leakier abstraction.

% ===========================================================================
\section{Case Study: A Longitudinal Single-User Deployment}
\label{sec:casestudy}

We report on one author's production workspace, which has accreted applications
through ordinary use over roughly two months. No host code was written for any of
these applications; each is a schema plus records, authored by the agent in
response to natural-language requests.

\paragraph{Scale.}
The workspace contains 26 collections holding $\sim$700 records. Their sizes vary
by three orders of magnitude (Table~\ref{tab:collections}), consistent with the
claim that the marginal cost of a new application is low enough to justify even a
single-record one.

\begin{table}[t]
\centering
\small
\begin{tabular}{@{}lr lr lr@{}}
\toprule
Collection & \#rec & Collection & \#rec & Collection & \#rec \\
\midrule
mag2-subscribers & 466 & engagements & 22 & recurring-bills & 6 \\
vocabulary       & 50  & restaurants & 20 & videos          & 6 \\
todos            & 23  & lessons-biology & 12 & worklog      & 6 \\
reading-list     & 10  & baseball-scout  & 11 & portfolio    & 4 \\
stock-quotes     & 10  & watchlist       & 10 & html-apps    & 4 \\
tax-payments     & 7   & clients         & 3  & \emph{(13 more)} & $\le 2$ \\
\bottomrule
\end{tabular}
\caption{Representative collections by record count in the studied workspace
(26 total). Sizes span from a single-record profile to 466 newsletter
subscribers.}
\label{tab:collections}
\end{table}

\paragraph{Mechanisms exercised in practice.}
The deployment uses every major mechanism described in \S\ref{sec:collections}:
relationships (an \code{invoice} references a \code{client} and embeds the issuer
profile), derived computation (invoice totals from line items and tax),
host-driven recurrence (monthly and annual obligations via \code{spawn}), and
completion/notification gating (todos, tax payments).

\paragraph{Bespoke views.}
Six collections carry custom HTML views authored by the agent on request: a
restaurant guide rendered as a transit-style map, a movie watchlist as a poster
gallery, a vocabulary list as a flashcard drill, a portfolio as an allocation
chart, a scouting collection as per-entity radar charts, and a subscriber
collection as a trend line. Each is software with an addressable audience of
one---software no vendor would build---produced for the cost of a request. The
owner never wrote or read the HTML.

\paragraph{Feeds.}
Three collections are feeds: an hourly newspaper-technology feed (64 records at
time of writing), and two daily podcast feeds (100 records each). One podcast feed
is read--write: the schema layers the owner's own rating, notes, and
resume-position fields on top of the ingested episode metadata, and the host
serves a custom gallery view over the combined record.

\paragraph{Observations.}
Two qualitative observations stand out. (i) \emph{No per-application host code} was
required for any of the 26 applications; the host's generic capabilities sufficed,
which is the load-bearing claim of applications-as-data. (ii) The applications
\emph{diverge sharply from any mass-market product}: a transit-map restaurant
guide and a read--write podcast tracker with personal resume positions reflect one
user's specific needs, supporting the view that a near-zero marginal build cost
shifts the natural unit of software from the mass-market product toward the
individual.

% ===========================================================================
\section{Discussion and Limitations}
\label{sec:limitations}

\paragraph{The escape-hatch hazard.}
The expressiveness that escape hatches restore is also unconstrained: a custom
view or a seeded-chat action runs outside the declarative guarantees. We mitigate
this for views with sandboxing and capability-scoped tokens (\S\ref{sec:collections}),
but the general tension---more expressiveness, fewer guarantees---is fundamental
and must be managed per escape, not eliminated.

\paragraph{Non-replayable execution.}
Because the operating step is a nondeterministic model call, a run cannot be
replayed byte-for-byte, and a self-modification's effect cannot be A/B-replayed
against the prior behavior on identical inputs. This is the deliberate trade
against log-primary designs~\citep{ynakajima2026log}. The mitigation is
configuration-level rather than run-level: a git-backed workspace makes every
schema and view edit auditable and revertable, which is the granularity that
matters for an agent editing its own applications.

\paragraph{Correctness has no second audience.}
Mass software accumulates correctness through many users and years; an application
with an audience of one does not. A wrong poster gallery is harmless, but a wrong
financial or tax view is not, and only the model and the single user ever inspect
it. The reliability dial is the relevant defense---push consequential computation
into host-guaranteed derived fields rather than leaving it to a generated
view---but the residual risk is real and is the price of building for one.

\paragraph{Threats to validity.}
Our empirical evidence is a single user, a single system, and an author-operated
deployment; it demonstrates feasibility and the no-host-code property but does not
establish generality, and it reports usage rather than a controlled comparison.
Quantifying authoring effort, error rates, and how the reliability dial should be
set across users and domains is future work.

% ===========================================================================
\section{Conclusion}
\label{sec:conclusion}

A domain-specific language is a powerful harness for an autonomous agent precisely
because it relinquishes expressiveness for a small, validated, inspectable
surface. Collections applies that principle to applications themselves: a schema
defines the application, record files are the database, the host runs the
deterministic half, and the model operates inside the schema. The result is
software an agent can author and refine on request, with the reliability of a
declarative contract and the flexibility of a language model held apart by a
deliberate dial. Building applications shifts from writing imperative code to
declaring---and validating---harnesses; the records are data, the schema is the
program, and the model is the runtime.

% ===========================================================================
\section*{Availability}
MulmoClaude and the Collections implementation are open source. The schema
validator, derived-formula evaluator, recurrence reconciler, custom-view sandbox,
and feed engine referenced above are part of that codebase, available at
\url{https://github.com/receptron/mulmoclaude}.

\bibliographystyle{plainnat}
\bibliography{refs}

\end{document}
